%==================================================================
% Ini adalah bab 3
% Silahkan edit sesuai kebutuhan, baik menambah atau mengurangi \section, \subsection
%==================================================================

\chapter[METODE PENELITIAN]{\\ METODE PENELITIAN}

\section{Jenis Penelitian}
\begin{sectioncontent}
    \hspace{\parindent}Penelitian ini menggunakan pendekatan eksperimen komparatif. Pendekatan ini dipilih karena penelitian membandingkan dua model \textit{deep learning}, yaitu YOLOv8 Nano dan YOLO11 Nano, pada tugas klasifikasi penyakit daun tomat. Perbandingan dilakukan berdasarkan dua aspek utama, yaitu akurasi klasifikasi dan efisiensi inferensi ketika model dijalankan pada perangkat \textit{edge}.

    Eksperimen dilakukan dengan menggunakan dataset, konfigurasi pelatihan, dan skenario pengujian yang sama untuk kedua model. Agar hasil perbandingan tidak hanya bergantung pada satu kali proses pelatihan, masing-masing model dilatih menggunakan beberapa variasi seed. Hasil dari setiap seed kemudian dibandingkan secara berpasangan, sehingga perbedaan performa kedua model dapat dianalisis dengan lebih hati-hati.

    Hasil dari penelitian ini digunakan untuk menentukan model yang lebih sesuai untuk diterapkan pada sistem klasifikasi penyakit daun tomat berbasis perangkat \textit{edge}. Penentuan tersebut tidak hanya didasarkan pada nilai akurasi, tetapi juga mempertimbangkan waktu inferensi, FPS, ukuran model, dan penggunaan memori pada perangkat pengujian.
\end{sectioncontent}

\section{Objek Penelitian}
\begin{sectioncontent}
    \hspace{\parindent}Objek dalam penelitian ini adalah citra daun tomat yang terdiri dari beberapa kelas penyakit dan satu kelas daun sehat. Citra tersebut digunakan sebagai data masukan untuk proses klasifikasi menggunakan model YOLOv8 Nano dan YOLO11 Nano. Penelitian ini berfokus pada tugas \textit{image classification}, sehingga keluaran model berupa label kelas dari keseluruhan citra, bukan lokasi objek atau area penyakit pada citra.

    \newpage

    Kelas yang digunakan dalam penelitian ini terdiri dari \textit{Early Blight}, \textit{Late Blight}, \textit{Septoria Leaf Spot}, \textit{Leaf Mold}, \textit{Yellow Leaf Curl Virus}, dan \textit{Healthy}. Penentuan kelas tersebut disesuaikan dengan struktur dataset publik yang digunakan dalam penelitian.
\end{sectioncontent}

\section{Tahapan Penelitian}
\begin{sectioncontent}
    \hspace{\parindent}Tahapan penelitian disusun agar proses eksperimen dapat dilakukan secara sistematis, mulai dari studi literatur, persiapan dataset, pelatihan model, hingga analisis hasil pengujian. Secara umum, tahapan penelitian ini adalah sebagai berikut:

    \begin{enumerate}[leftmargin=0.75cm,label=\arabic*.]
        \item Melakukan studi literatur mengenai penyakit daun tomat, klasifikasi citra, model YOLO, dan perangkat \textit{edge}.
        \item Mengumpulkan dataset citra penyakit daun tomat dari sumber publik.
        \item Memeriksa struktur dataset dan jumlah data pada setiap kelas.
        \item Memeriksa pembagian dataset awal yang telah tersedia dalam struktur data latih, data validasi, dan data uji.
        \item Melakukan \textit{preprocessing} dataset agar sesuai dengan kebutuhan model klasifikasi YOLO.
        \item Melatih model YOLOv8 Nano dan YOLO11 Nano menggunakan konfigurasi yang sama pada beberapa variasi seed.
        \item Mengevaluasi performa klasifikasi model menggunakan data uji yang sama untuk setiap model dan setiap seed.
        \item Mengekspor model ke format yang sesuai untuk pengujian pada perangkat \textit{edge}.
        \item Menjalankan pengujian inferensi model pada Raspberry Pi 5.
        \item Menganalisis perbandingan akurasi dan efisiensi kedua model secara deskriptif dan menggunakan uji Wilcoxon Signed-Rank.
        \item Menarik kesimpulan berdasarkan hasil pengujian yang diperoleh.
    \end{enumerate}
\end{sectioncontent}

\section{Dataset Penelitian}
\begin{sectioncontent}
    \hspace{\parindent}Dataset yang digunakan dalam penelitian ini adalah \textit{Tomato Leaf Disease Classification Dataset in Pakistan} \citep{malik2026dataset} yang tersedia secara publik. Dataset ini terdiri dari 7.200 citra daun tomat yang dikumpulkan dari kondisi pertanian lapangan di Pakistan. Citra pada dataset tersebut mencakup enam kelas, yaitu \textit{Early Blight}, \textit{Late Blight}, \textit{Septoria Leaf Spot}, \textit{Leaf Mold}, \textit{Yellow Leaf Curl Virus}, dan \textit{Healthy}.

    Dataset yang digunakan dalam penelitian ini telah tersedia dalam struktur data latih, data validasi, dan data uji. Oleh karena itu, penelitian ini mengikuti struktur pembagian dataset awal tanpa melakukan pembagian ulang. Penggunaan struktur awal tersebut dipilih agar eksperimen tetap konsisten dengan susunan dataset yang telah tersedia.

    Struktur direktori dataset disesuaikan dengan format klasifikasi citra pada Ultralytics YOLO, yaitu terdiri dari folder \textit{train}, \textit{val}, dan \textit{test}. Setiap folder memiliki subdirektori berdasarkan nama kelas. Data latih digunakan untuk proses pelatihan model, data validasi digunakan untuk memantau performa model selama pelatihan, sedangkan data uji digunakan untuk mengevaluasi performa akhir model.

    Penelitian ini tidak menentukan seed untuk pembagian dataset karena struktur \textit{train}, \textit{val}, dan \textit{test} telah tersedia dari dataset awal. Seed digunakan pada proses pelatihan model, bukan untuk membagi ulang dataset. Dengan demikian, setiap model tetap diuji pada data uji yang sama, sedangkan variasi seed digunakan untuk melihat kestabilan hasil pelatihan.

    Dataset disusun ke dalam format direktori yang sesuai dengan kebutuhan tugas klasifikasi pada Ultralytics YOLO. Struktur dataset terdiri dari tiga bagian utama, yaitu data latih, data validasi, dan data uji. Setiap bagian memiliki subdirektori berdasarkan nama kelas.

    \begin{verbatim}
tomato-leaf-disease/
|-- train/
|   |-- Tomato_Early_blight_leaf/
|   |-- Tomato_Healthy_leaf/
|   |-- Tomato_leaf_late_blight/
|   |-- Tomato_leaf_yellow_curl_virus/
|   |-- Tomato_mold_leaf/
|   |-- Tomato_Septoria_leaf_spot/
|
|-- val/
|   |-- Tomato_Early_blight_leaf/
|   |-- Tomato_Healthy_leaf/
|   |-- Tomato_leaf_late_blight/
|   |-- Tomato_leaf_yellow_curl_virus/
|   |-- Tomato_mold_leaf/
|   |-- Tomato_Septoria_leaf_spot/
|
|-- test/
    |-- Tomato_Early_blight_leaf/
    |-- Tomato_Healthy_leaf/
    |-- Tomato_leaf_late_blight/
    |-- Tomato_leaf_yellow_curl_virus/
    |-- Tomato_mold_leaf/
    |-- Tomato_Septoria_leaf_spot/
    \end{verbatim}

    Pembagian dataset pada sumber data telah disusun dengan proporsi yang seimbang pada setiap kelas. Dalam penelitian ini, pembagian tersebut digunakan sebagaimana struktur awal dataset tanpa proses pembagian ulang.
\end{sectioncontent}

\section{Perangkat Keras dan Perangkat Lunak}
\begin{sectioncontent}
    \hspace{\parindent}Penelitian ini menggunakan perangkat keras dan perangkat lunak untuk mendukung proses pelatihan, pengujian, dan implementasi model pada perangkat \textit{edge}. Proses pelatihan dilakukan pada komputer dengan dukungan GPU agar pelatihan model dapat berjalan lebih cepat. Sementara itu, pengujian efisiensi inferensi dilakukan pada Raspberry Pi 5 sebagai representasi perangkat \textit{edge}.

    \subsection{Perangkat Keras}
    \begin{subsectioncontent}
        \hspace{\parindent}Perangkat keras yang digunakan dalam penelitian ini adalah sebagai berikut:

        \begin{enumerate}[leftmargin=0.75cm,label=\arabic*.]
            \item \textit{Desktop PC} dengan spesifikasi prosesor AMD Ryzen 5, RAM 16 GB, dan GPU NVIDIA RTX 4060 untuk pelatihan model.
            \item Raspberry Pi 5 dengan RAM 8 GB sebagai perangkat \textit{edge} untuk pengujian inferensi.
            \item Media penyimpanan untuk menyimpan dataset, hasil pelatihan, dan model hasil ekspor.
        \end{enumerate}
    \end{subsectioncontent}

    \subsection{Perangkat Lunak}
    \begin{subsectioncontent}
        \hspace{\parindent}Perangkat lunak yang digunakan dalam penelitian ini adalah sebagai berikut:

        \begin{enumerate}[leftmargin=0.75cm,label=\arabic*.]
            \item Python sebagai bahasa pemrograman utama.
            \item Ultralytics YOLO untuk pelatihan dan evaluasi model YOLOv8 Nano dan YOLO11 Nano.
            \item PyTorch sebagai \textit{backend} pelatihan model.
            \item CUDA untuk akselerasi pelatihan menggunakan GPU NVIDIA.
            \item ONNX Runtime untuk menjalankan inferensi model pada perangkat \textit{edge}.
            \item NumPy, Pandas, dan Scikit-learn untuk pengolahan data serta evaluasi metrik.
            \item SciPy untuk melakukan uji Wilcoxon Signed-Rank.
            \newpage
            \item Sistem operasi Windows pada perangkat pelatihan dan sistem operasi berbasis Linux pada Raspberry Pi 5.
        \end{enumerate}
    \end{subsectioncontent}
\end{sectioncontent}

\section{\textit{Preprocessing} Data}
\begin{sectioncontent}
    \hspace{\parindent}\textit{Preprocessing} dilakukan untuk memastikan citra dapat digunakan pada proses pelatihan, evaluasi, dan inferensi model. Pada tahap pelatihan, preprocessing mengikuti mekanisme yang digunakan oleh Ultralytics YOLO untuk tugas klasifikasi citra. Citra diproses dengan ukuran masukan 224 x 224 piksel sesuai dengan konfigurasi pelatihan yang digunakan.

    Tahapan \textit{preprocessing} dalam penelitian ini meliputi pemeriksaan struktur folder dataset, pemeriksaan jumlah data pada setiap kelas, penyesuaian ukuran citra menjadi 224 x 224 piksel, serta normalisasi nilai piksel sesuai kebutuhan model. Dataset yang digunakan sudah memiliki struktur \textit{train}, \textit{val}, dan \textit{test}, sehingga penelitian ini tidak melakukan pembagian ulang dataset.

    Pada tahap pengujian inferensi menggunakan model ONNX, preprocessing dilakukan secara manual melalui script benchmark. Citra dibaca dalam format RGB, diubah ukurannya menjadi 224 x 224 piksel, dinormalisasi ke rentang nilai 0 sampai 1, kemudian diubah dari format HWC menjadi CHW dan ditambahkan dimensi batch. Tahapan tersebut dilakukan agar bentuk input sesuai dengan kebutuhan model YOLO versi klasifikasi dalam format ONNX.
\end{sectioncontent}

\section{Model yang Digunakan}
\begin{sectioncontent}
    \hspace{\parindent}Model yang digunakan dalam penelitian ini adalah YOLOv8 Nano dan YOLO11 Nano versi klasifikasi. Kedua model dipilih karena termasuk varian ringan dari keluarga YOLO, sehingga relevan untuk diuji pada perangkat dengan sumber daya komputasi terbatas.

    \subsection{YOLOv8 Nano}
    \begin{subsectioncontent}
        \hspace{\parindent}YOLOv8 Nano merupakan varian ringan dari YOLOv8 yang dirancang untuk memiliki ukuran model kecil dan waktu inferensi yang cepat. Dalam penelitian ini, YOLOv8 Nano digunakan sebagai salah satu model pembanding untuk mengklasifikasikan penyakit daun tomat berbasis citra.
    \end{subsectioncontent}

    \subsection{YOLO11 Nano}
    \begin{subsectioncontent}
        \hspace{\parindent}YOLO11 Nano merupakan varian ringan dari YOLO11 yang digunakan sebagai model pembanding terhadap YOLOv8 Nano. Model ini dipilih karena berasal dari generasi YOLO yang lebih baru dan tetap diarahkan pada keseimbangan antara performa dan efisiensi. Perbandingan antara YOLOv8 Nano dan YOLO11 Nano dilakukan untuk mengetahui perbedaan performa keduanya pada tugas klasifikasi penyakit daun tomat.
    \end{subsectioncontent}
\end{sectioncontent}

\section{Skenario Pelatihan Model}
\begin{sectioncontent}
    \hspace{\parindent}Pelatihan model dilakukan menggunakan YOLOv8 Nano dan YOLO11 Nano versi klasifikasi. Kedua model dilatih menggunakan dataset, ukuran citra, jumlah epoch, batch size, dan perangkat pelatihan yang sama. Penggunaan konfigurasi yang sama bertujuan agar perbandingan hasil lebih diarahkan pada perbedaan karakteristik model, bukan pada perbedaan skenario pelatihan.

    Untuk mengurangi ketergantungan hasil pada satu kali proses pelatihan, masing-masing model dilatih sebanyak 20 kali menggunakan seed 1 sampai 20. Setiap seed pada YOLOv8 Nano dipasangkan dengan seed yang sama pada YOLO11 Nano. Dengan cara tersebut, hasil pengujian dapat dianalisis sebagai data berpasangan pada tahap uji statistik. Konfigurasi pelatihan yang digunakan ditampilkan pada Tabel \ref{tab:konfigurasi-pelatihan}.

    \begin{table}[H]
        \centering
        \caption{Konfigurasi Pelatihan Model}
        \label{tab:konfigurasi-pelatihan}
        \begin{tabular}{|l|l|}
            \hline
            \textbf{Parameter} & \textbf{Nilai} \\ \hline
                Model 1 & YOLOv8 Nano Classification \\ \hline
                Model 2 & YOLO11 Nano Classification \\ \hline
                Ukuran citra & 224 x 224 piksel \\ \hline
                Jumlah epoch & 50 \\ \hline
                Batch size & 16 \\ \hline
                Workers & 2 \\ \hline
                Cache & False \\ \hline
                Seed & 1 sampai 20 \\ \hline
                Perangkat pelatihan & Desktop PC dengan GPU NVIDIA RTX 4060 \\ \hline
        \end{tabular}
    \end{table}

    Perintah pelatihan yang digunakan adalah sebagai berikut:

    \begin{verbatim}
    yolo classify train model=yolov8n-cls.pt \
    data=tomato-leaf-disease epochs=50 imgsz=224 batch=16 \
    device=0 workers=2 cache=False seed=<nilai_seed> \
    deterministic=True pretrained=True project=runs/yolov8n \
    name=seed-<nilai_seed> exist_ok=True

    yolo classify train model=yolo11n-cls.pt \
    data=tomato-leaf-disease epochs=50 imgsz=224 batch=16 \
    device=0 workers=2 cache=False seed=<nilai_seed> \
    deterministic=True pretrained=True project=runs/yolo11n \
    name=seed-<nilai_seed> exist_ok=True
    \end{verbatim}

    Kedua model dilatih dengan konfigurasi yang sama agar hasil perbandingan lebih konsisten. Apabila terdapat perbedaan performa pada tahap evaluasi, perbedaan tersebut dapat dianalisis sebagai hasil dari karakteristik masing-masing model, dengan tetap mempertimbangkan variasi hasil akibat perbedaan seed.
\end{sectioncontent}

\section{Skenario Pengujian Model}
\begin{sectioncontent}
    \hspace{\parindent}Pengujian model dilakukan dalam dua tahap, yaitu pengujian performa klasifikasi dan pengujian efisiensi pada perangkat \textit{edge}. Pengujian performa klasifikasi digunakan untuk mengetahui kemampuan model dalam mengklasifikasikan citra daun tomat ke kelas yang benar. Sementara itu, pengujian efisiensi digunakan untuk mengetahui kecepatan dan kelayakan model ketika dijalankan pada Raspberry Pi 5.

    \subsection{Pengujian Klasifikasi}
    \begin{subsectioncontent}
        \hspace{\parindent}Pengujian klasifikasi dilakukan menggunakan data uji yang tidak digunakan dalam proses pelatihan. Setiap model hasil pelatihan diuji pada data uji yang sama, sehingga terdapat 20 hasil pengujian untuk YOLOv8 Nano dan 20 hasil pengujian untuk YOLO11 Nano. Hasil prediksi model dibandingkan dengan label sebenarnya untuk menghitung metrik evaluasi, seperti akurasi, \textit{precision}, \textit{recall}, dan \textit{F1-score}. Selain itu, \textit{confusion matrix} digunakan untuk melihat pola kesalahan klasifikasi, terutama pada kelas-kelas yang memiliki gejala visual mirip.
    \end{subsectioncontent}

    \subsection{Pengujian pada Perangkat \textit{Edge}}
    \begin{subsectioncontent}
        \hspace{\parindent}Pengujian pada perangkat \textit{edge} dilakukan dengan menjalankan model hasil pelatihan pada Raspberry Pi 5. Model yang telah dilatih diekspor ke format ONNX agar dapat dijalankan menggunakan ONNX Runtime pada perangkat \textit{edge}. Pengujian ini mengukur efisiensi inferensi kedua model pada lingkungan komputasi terbatas.

        Parameter yang diamati dalam pengujian pada perangkat \textit{edge} meliputi:

        \begin{enumerate}[leftmargin=0.75cm,label=\arabic*.]
            \item Waktu inferensi rata-rata per citra.
            \item Waktu total \textit{pipeline} per citra.
            \item Nilai FPS (\textit{Frame Per Second}).
            \item Penggunaan memori saat inferensi.
            \item Ukuran file model hasil ekspor.
        \end{enumerate}

        Pengujian dilakukan dengan ukuran \textit{batch} 1 menggunakan ONNX Runtime dengan CPUExecutionProvider. Sebelum pengukuran utama, dilakukan proses warm-up sebanyak 10 citra untuk membantu menstabilkan waktu inferensi. Script benchmark mencatat waktu \textit{preprocessing}, waktu inferensi, waktu \textit{postprocessing}, dan waktu total \textit{pipeline}. Metrik utama yang digunakan dalam perbandingan efisiensi adalah waktu inferensi rata-rata, waktu total \textit{pipeline}, dan FPS, karena metrik tersebut menggambarkan kecepatan model ketika dijalankan pada Raspberry Pi 5.

        Penggunaan memori diamati menggunakan nilai \textit{Resident Set Size} (RSS) dari proses Python selama benchmark berjalan. Nilai memori puncak yang dilaporkan merupakan RSS puncak yang tersampling selama proses pengujian, bukan pengukuran memori perangkat secara keseluruhan.
    \end{subsectioncontent}
\end{sectioncontent}

\section{Metrik Evaluasi}
\begin{sectioncontent}
    \hspace{\parindent}Metrik evaluasi digunakan untuk mengukur performa model dari sisi akurasi klasifikasi dan efisiensi inferensi. Pada penelitian ini, metrik evaluasi yang digunakan terdiri dari akurasi, \textit{precision}, \textit{recall}, \textit{F1-score}, waktu inferensi, FPS, penggunaan memori, dan ukuran model. Karena penelitian ini menggunakan klasifikasi multikelas, nilai \textit{precision}, \textit{recall}, dan \textit{F1-score} juga dihitung dalam bentuk \textit{macro average} agar performa setiap kelas dapat diperhitungkan secara seimbang.

    \subsection{Akurasi}
    \begin{subsectioncontent}
        \hspace{\parindent}Akurasi menunjukkan perbandingan antara jumlah prediksi yang benar dengan seluruh data yang diuji. Dalam klasifikasi multikelas, akurasi digunakan untuk mengetahui persentase keseluruhan citra yang berhasil diklasifikasikan dengan benar oleh model. Rumus akurasi adalah sebagai berikut:

        \begin{equation}
            Accuracy = \frac{\text{Jumlah Prediksi Benar}}{\text{Jumlah Seluruh Data Uji}}
        \end{equation}
    \end{subsectioncontent}

    \subsection{\textit{Precision}}
    \begin{subsectioncontent}
        \hspace{\parindent}\textit{Precision} menunjukkan tingkat ketepatan model ketika memprediksi suatu kelas. Nilai ini menggambarkan seberapa banyak prediksi pada suatu kelas yang benar-benar sesuai dengan kelas sebenarnya. Rumus \textit{precision} adalah sebagai berikut:

        \begin{equation}
            Precision = \frac{TP}{TP + FP}
        \end{equation}
    \end{subsectioncontent}

    \subsection{\textit{Recall}}
    \begin{subsectioncontent}
        \hspace{\parindent}\textit{Recall} menunjukkan kemampuan model dalam menemukan seluruh data yang seharusnya termasuk ke dalam suatu kelas. Nilai ini digunakan untuk melihat seberapa baik model mengenali data dari kelas tertentu. Rumus \textit{recall} adalah sebagai berikut:

        \begin{equation}
            Recall = \frac{TP}{TP + FN}
        \end{equation}
    \end{subsectioncontent}

    \subsection{\textit{F1-score}}
    \begin{subsectioncontent}
        \hspace{\parindent}\textit{F1-score} merupakan rata-rata harmonik dari \textit{precision} dan \textit{recall}. Metrik ini digunakan untuk memberikan gambaran keseimbangan antara ketepatan prediksi dan kemampuan model dalam menemukan data pada suatu kelas. Rumus \textit{F1-score} adalah sebagai berikut:

        \begin{equation}
            F1\text{-}score = 2 \times \frac{Precision \times Recall}{Precision + Recall}
        \end{equation}
    \end{subsectioncontent}

    \subsection{\textit{Macro Average}}
    \begin{subsectioncontent}
        \hspace{\parindent}Pada klasifikasi multikelas, nilai \textit{precision}, \textit{recall}, dan \textit{F1-score} dapat dihitung untuk setiap kelas. Agar performa seluruh kelas dapat dilihat secara seimbang, penelitian ini menggunakan pendekatan \textit{macro average}. Pendekatan ini menghitung nilai rata-rata dari setiap kelas tanpa memberikan bobot lebih besar pada kelas tertentu.

        \begin{equation}
            \operatorname{Macro\ Precision} = \frac{1}{K} \sum_{i=1}^{K} \operatorname{Precision}_{i}
        \end{equation}

        \begin{equation}
            \operatorname{Macro\ Recall} = \frac{1}{K} \sum_{i=1}^{K} \operatorname{Recall}_{i}
        \end{equation}

        \begin{equation}
            \operatorname{Macro\ F1\text{-}score} = \frac{1}{K} \sum_{i=1}^{K} F1_{i}
        \end{equation}

        Keterangan $K$ menunjukkan jumlah kelas yang digunakan dalam penelitian. Dengan menggunakan \textit{macro average}, setiap kelas memiliki pengaruh yang sama terhadap nilai akhir, sehingga hasil evaluasi tidak hanya mencerminkan performa pada kelas yang jumlah datanya lebih besar atau lebih mudah dikenali.
    \end{subsectioncontent}

    \subsection{Waktu Inferensi}
    \begin{subsectioncontent}
        \hspace{\parindent}Waktu inferensi menunjukkan waktu yang dibutuhkan model untuk memproses satu citra dan menghasilkan prediksi. Semakin rendah waktu inferensi, semakin cepat model dalam melakukan klasifikasi.
    \end{subsectioncontent}

    \subsection{\textit{Frame Per Second}}
    \begin{subsectioncontent}
        \hspace{\parindent}\textit{Frame Per Second} atau FPS menunjukkan jumlah citra yang dapat diproses oleh model dalam satu detik. FPS dihitung berdasarkan waktu rata-rata menggunakan rumus berikut:

        \begin{equation}
            FPS = \frac{1000}{\text{Average Time (ms)}}
        \end{equation}

        Dalam penelitian ini, FPS dapat dihitung berdasarkan waktu inferensi saja maupun berdasarkan waktu total \textit{pipeline}. Semakin tinggi nilai FPS, semakin baik kemampuan model dalam memproses citra secara cepat pada perangkat \textit{edge}.
    \end{subsectioncontent}

    \subsection{Penggunaan Memori dan Ukuran Model}
    \begin{subsectioncontent}
        \hspace{\parindent}Penggunaan memori digunakan untuk melihat kebutuhan sumber daya model ketika dijalankan pada Raspberry Pi 5. Nilai yang diamati adalah \textit{Resident Set Size} (RSS) puncak yang tersampling selama proses pengujian. Sementara itu, ukuran model menunjukkan besar file model hasil ekspor dalam format ONNX. Kedua metrik ini digunakan sebagai pertimbangan tambahan karena perangkat \textit{edge} umumnya memiliki keterbatasan penyimpanan dan memori.
    \end{subsectioncontent}
\end{sectioncontent}

\section{Teknik Analisis Data}
\begin{sectioncontent}
    \hspace{\parindent}Analisis data dilakukan dengan membandingkan hasil pengujian YOLOv8 Nano dan YOLO11 Nano berdasarkan metrik yang telah ditentukan. Hasil pengujian klasifikasi dianalisis menggunakan akurasi, \textit{macro precision}, \textit{macro recall}, \textit{macro F1-score}, hasil evaluasi per kelas, dan \textit{confusion matrix}. Sementara itu, hasil pengujian efisiensi pada perangkat \textit{edge} dianalisis berdasarkan waktu inferensi, waktu total \textit{pipeline}, FPS, penggunaan memori, dan ukuran model.

    Setiap hasil pengujian YOLOv8 Nano dipasangkan dengan hasil pengujian YOLO11 Nano pada seed yang sama. Pasangan data tersebut kemudian dianalisis menggunakan uji Wilcoxon Signed-Rank untuk melihat apakah perbedaan hasil kedua model signifikan secara statistik. Uji ini digunakan karena data yang dibandingkan bersifat berpasangan dan jumlah pasangan pengujian relatif terbatas.

    Taraf signifikansi yang digunakan dalam penelitian ini adalah $\alpha = 0{,}05$. Apabila nilai \textit{p-value} lebih kecil dari 0,05, maka perbedaan kedua model pada metrik tersebut dinyatakan signifikan secara statistik. Sebaliknya, apabila nilai \textit{p-value} lebih besar atau sama dengan 0,05, maka perbedaan kedua model tidak dinyatakan signifikan. Ukuran model tetap dilaporkan sebagai hasil deskriptif karena nilainya merupakan properti file model hasil ekspor dan tidak berubah antar seed.

    Hasil pengujian disajikan dalam bentuk tabel dan grafik agar perbandingan kedua model lebih mudah dianalisis. Model yang memiliki akurasi lebih tinggi, waktu inferensi lebih rendah, dan ukuran model lebih kecil dinilai lebih sesuai untuk diterapkan pada sistem klasifikasi penyakit daun tomat berbasis perangkat \textit{edge}. Namun, apabila terdapat perbedaan antara akurasi dan efisiensi, analisis dilakukan dengan mempertimbangkan keseimbangan antara kedua aspek tersebut.
\end{sectioncontent}

\section{Rancangan Hasil Pengujian}
\begin{sectioncontent}
    \hspace{\parindent}Rancangan hasil pengujian digunakan untuk menggambarkan bentuk data yang dianalisis pada bab hasil dan pembahasan. Perbandingan hasil pengujian kedua model disajikan dalam tabel berikut:

    \begin{table}[H]
        \centering
        \caption{Rancangan Tabel Perbandingan Akurasi Model}
        \begin{tabularx}{\linewidth}{|l|>{\centering\arraybackslash}X|>{\centering\arraybackslash}X|>{\centering\arraybackslash}X|>{\centering\arraybackslash}X|}
            \hline
            \textbf{Model} & \textbf{Akurasi} & \textbf{\shortstack{Macro\\Precision}} & \textbf{\shortstack{Macro\\Recall}} & \textbf{\shortstack{Macro\\F1-score}} \\
            \hline
            YOLOv8 Nano & - & - & - & - \\
            \hline
            YOLO11 Nano & - & - & - & - \\
            \hline
        \end{tabularx}
    \end{table}

    \begin{table}[H]
        \centering
        \caption{Rancangan Tabel Perbandingan Efisiensi pada Perangkat Edge}
        \begin{tabularx}{\linewidth}{|l|>{\centering\arraybackslash}X|>{\centering\arraybackslash}X|>{\centering\arraybackslash}X|>{\centering\arraybackslash}X|}
            \hline
            \textbf{Model} & \textbf{\shortstack{Waktu\\Inferensi}} & \textbf{FPS} & \textbf{\shortstack{Memori\\Puncak\\Teramati}} & \textbf{\shortstack{Ukuran\\Model}} \\
            \hline
            YOLOv8 Nano & - & - & - & - \\
            \hline
            YOLO11 Nano & - & - & - & - \\
            \hline
        \end{tabularx}
    \end{table}

    \begin{table}[H]
        \centering
        \caption{Rancangan Tabel Hasil Uji Wilcoxon Signed-Rank}
        \begin{tabularx}{\linewidth}{|l|>{\centering\arraybackslash}X|>{\centering\arraybackslash}X|>{\centering\arraybackslash}X|}
            \hline
            \textbf{Metrik} & \textbf{Statistik Uji} & \textbf{\textit{p-value}} & \textbf{Keputusan} \\
            \hline
            - & - & - & - \\
            \hline
        \end{tabularx}
    \end{table}

    Tabel tersebut digunakan sebagai dasar untuk membandingkan performa kedua model. Penentuan model yang lebih sesuai tidak hanya didasarkan pada akurasi tertinggi, tetapi juga mempertimbangkan efisiensi inferensi dan hasil uji statistik ketika model dijalankan pada perangkat \textit{edge}.
\end{sectioncontent}
